\section{Strings}

\subsection{Hashing}
\algoinfo{|s|}{|s|}{Calcula el hash de los prefijos de un string. Admite concatenacion y potencias.}
\inputcode{cpp}{latex_src/strings/hashing.cpp}{hashing.cpp}

\subsection{KMP}
Complejidad: Tiempo $O(|s|)$ - Memoria extra $O(|s|)$.
\begin{minted}{cpp}
vi prefix_function(const string &s){
    int n = sz(s);
    vi pi(n);
    for(int i = 1; i < n; ++i){
        int j = pi[i - 1];
        while(j && s[i] != s[j]) j = pi[j - 1];
        pi[i] = j + (s[i] == s[j]);
    } return pi;
}
\end{minted}



\textbf{Autómata de KMP.} Complejidad: Tiempo $O(|s|k)$ - Memoria extra $O(|s|k)$, donde $k$ es el tamaño del alfabeto.
\begin{minted}{cpp}
void compute_automaton(const string &s, vvi& aut){
    s += '#';
    int n = sz(s);
    vi pi = prefix_function(s);
    aut.assign(n, vi(26));
    for(int i = 0; i < n; ++i){
        for(int c = 0; c < 26; ++c){
            if(i && 'a' + c != s[i])
                aut[i][c] = aut[pi[i - 1]][c];
            else aut[i][c] = i + ('a'+c == s[i]);
        }
    }
}
\end{minted}



\subsection{Suffix Automata}
Complejidad: Tiempo $O(|s|)$, memoria $O(|S|)$. Construye un autómata sobre los sufijos de $s$, donde cada nodo terminal es un sufijo de $s$. Esta construcción se hace online. Donde $alpha$ es el tamaño del alfabeto y $offset$ el primer caracter. Si se hace uso de los nodos terminales llamar a markTerminalNodes() primero (este método no es online).
\begin{minted}{cpp}
constexpr short alpha = 26;
constexpr char offset = 'a';
struct state{
    //max(state) = len, min(state) = sa[link].len
    int len, link;
    bool is_terminal;
    array<int,alpha> next;
    state(){
        len = 0; link = -1; is_terminal = false;
        next.fill(0);
    }
};
struct SuffixAutomaton{
    int n;
    vector<state> sa;
    int sz = 1, last = 0;
    SuffixAutomaton(int n) : n(n), sa(2*n + 1){}
    void add(char ch_){
        short ch = ch_-offset;
        int curr = sz++;
        sa[curr].len = sa[last].len + 1;
        int p = last;
        while(p!=-1 && !sa[p].next[ch]){
            sa[p].next[ch] = curr;
            p = sa[p].link;
        }
        if(p == -1){
            sa[curr].link = 0;
        }else{
            int q = sa[p].next[ch];
            if(sa[p].len+1 == sa[q].len){
                sa[curr].link = q;
            }else{
                int clone = sz++;
                sa[clone].len = sa[p].len + 1;
                sa[clone].next = sa[q].next;
                sa[clone].link = sa[q].link;
                while(p!=-1 && sa[p].next[ch] == q){
                    sa[p].next[ch] = clone;
                    p = sa[p].link;
                }
                sa[q].link = sa[curr].link = clone;
            }
        }
        last = curr;
        return;
    }
    void markTerminalNodes(){
        int curr = last;
        while(curr){
            sa[curr].is_terminal = true;
            curr = sa[curr].link;
        }
        return;
    }
};
\end{minted}



\subsection{Suffix array}
\textbf{Construcción.} Complejidad: Tiempo $O(|s|\log(|s|))$ - Memoria $O(|s|)$. Calcula la permutación que corresponde a los sufijos ordenados lexicográficamente. \mintinline{cpp}{SA[i]} es el índice en el cual empieza el $i$-ésimo sufijo ordenado.
\begin{minted}{cpp}
int SA[MAXN], mrank[MAXN];
int tmpSA[MAXN], tmp_mrank[MAXN];
void counting_sort(int k, int n){
    int freqs[MAXN] = {};
    for(int i = 0; i < n; ++i){
        if(i + k < n) freqs[ mrank[i + k] ]++;
        else freqs[0]++;
    }
    int m = max(100, n);
    for(int i = 0, sfs = 0; i < m; ++i){
        int f = freqs[i];
        freqs[i] = sfs;
        sfs += f;
    }
    for(int i = 0; i < n; ++i){
        if(SA[i] + k < n) tmpSA[ freqs[mrank[ SA[i] + k ]]++ ] = SA[i];
        else tmpSA[ freqs[0]++ ] = SA[i];
    }
    for(int i = 0; i < n; ++i) SA[i] = tmpSA[i];
}
void buildSA(string &str){
    int n = sz(str);
    for(int i = 0; i < n; ++i){
        mrank[i] = str[i] - '#';
        SA[i] = i;
    }
    for(int k = 1; k < n; k <<= 1){
        counting_sort(k, n);
        counting_sort(0, n);
        int r = 0;
        tmp_mrank[ SA[0] ] = 0;
        for(int i = 1; i < n; ++i){
            if(mrank[ SA[i] ] != mrank[ SA[i - 1] ] || mrank[ SA[i] + k ] != mrank[ SA[i - 1] + k ])
                tmp_mrank[ SA[i] ] = ++r;
            else tmp_mrank[ SA[i] ] = r;
        }
        for(int i = 0; i < n; ++i) mrank[i] = tmp_mrank[i];
    }
}
inline bool suff_compare1(int idx,const string &pattern) {
    return (s.substr(idx).compare(0, sz(pattern), pattern) < 0);
}
inline bool suff_compare2(const string &pattern,int idx) {
    return (s.substr(idx).compare(0, sz(pattern), pattern) > 0);
}
pair<int,int> match(const string &pattern) {
    int *low = lower_bound (SA, SA + sz(s), pattern, suff_compare1);
    int *up = upper_bound (SA, SA + sz(s), pattern, suff_compare2);
    return make_pair((int)(low - SA),(int)(up - SA));
}
\end{minted}



\textbf{Prefijo común más largo.} Complejidad: Tiempo $O(|s|)$ - Memoria $O(|s|)$. Calcula la longitud del prefijo común más largo entre dos sufijos consecutivos (lexicográficamente) de \mintinline{cpp}{s}. $lcp[i]$ guarda la respuesta para el $i$-ésimo sufijo y el $(i - 1)$-ésimo sufijo.
\begin{minted}{cpp}
int lcp[MAXN];
void build_lcp(string &str){
    int n = sz(str);
    int phi[n];
    phi[SA[0]] = -1;
    for(int i = 1; i < n; ++i) phi[ SA[i] ] = SA[i - 1];
    int plcp[n];
    int k = 0;
    for(int i = 0; i < n; ++i){
        if(phi[i] == -1){
            plcp[i] = 0;
            continue;
        }
        while(i + k < n && phi[i] + k < n && str[i + k] == str[phi[i] + k]) k++;
        plcp[i] = k;
        k = max(k - 1, 0);
    }
    for(int i = 0; i < n; ++i) lcp[i] = plcp[SA[i]];
}
\end{minted}



\subsection{Aho-Corasick}
Construción en $O(mk)$, donde $m$ es el tamaño total de los strings y $k$ el tamaño del alfabeto.
\begin{minted}{cpp}
struct aho_corasick{
    const static int K = 26;
    const char index = 'a';
    struct vertex{
        int next[K];
        bool terminal = false;
        int p = -1;
        char p_edge;
        int link = -1;
        int terminal_link = -1;
        int go[K];
        vertex(int p = -1, char c = '$') : p(p), p_edge(c){
            memset(next, -1, K*sizeof(int));
            memset(go, -1, K*sizeof(int));
        }
    };
    vector<vertex> aho;
    aho_corasick(){ aho.resize(1); }
    void add_string(const string &s){
        int u = 0;
        for(char c : s){
            int e = c - index;
            if(aho[u].next[e] == -1){
                aho[u].next[e] = sz(aho);
                aho.emplace_back(u, c);
            }
            u = aho[u].next[e];
        } aho[u].terminal = 1;
    }
    int get_link(int u){
        if(aho[u].link == -1){
            if(!u || !aho[u].p) aho[u].link = 0;
            else aho[u].link = go(get_link(aho[u].p), aho[u].p_edge);
        } return aho[u].link;
    }
    int go(int u, char c){
        int e = c - index;
        if(aho[u].go[e] == -1){
            if(aho[u].next[e] != -1)
                aho[u].go[e] = aho[u].next[e];
            else aho[u].go[e] = !u ? 0 : go(get_link(u), c);
        } return aho[u].go[e];
    }
    int get_terminal_link(int u){
        if(aho[u].terminal_link == -1){
            if(!u || !aho[u].p)
                aho[u].terminal_link = 0;
            else aho[u].terminal_link = get_terminal_link(get_link(u));
        } return aho[u].terminal_link;
    }
};
\end{minted}



\subsection{Suffix tree}
\begin{minted}{cpp}
/// MEJORAR ESTA COSA, SOLO LO COPIE Y PEGUÉ porcuestionesdetiempo
const int inf = 1e9;
const int maxn = 1e6;
int s[maxn];
map<int, int> to[maxn];
//Root is the vertex 0
//f_pos[i] is the initial index with the letter of the edge that goes from the parent of i to i
//len[i] is the number of letters in the edge that enters in i
//slink[i] is the suffix link
int len[maxn], f_pos[maxn], slink[maxn];
int node, pos;
int sz = 1, n = 0;
int make_node(int _pos, int _len){
    f_pos[sz] = _pos;
    len [sz] = _len;
    return sz++;
}
void go_edge(){
    while(pos > len[to[node][s[n - pos]]]){
        node = to[node][s[n - pos]];
        pos -= len[node];
    }
}
void add_letter(int c){
    s[n++] = c;
    pos++;
    int last = 0;
    while(pos > 0){
        go_edge();
        int edge = s[n - pos];
        int &v = to[node][edge];
        int t = s[f_pos[v] + pos - 1];
        if(v == 0){
            v = make_node(n - pos, inf);
            //v = make_node(n - pos, 1);
            slink[last] = node;
            last = 0;
        } else if(t == c) {
            slink[last] = node;
            return;
        } else {
            int u = make_node(f_pos[v], pos - 1);
            to[u][c] = make_node(n - 1, inf);
            to[u][t] = v;
            f_pos[v] += pos - 1;
            len [v] -= pos - 1;
            v = u;
            slink[last] = u;
            last = u;
        }
        if(node == 0) pos--;
        else node = slink[node];
    }
}
void correct(int s_size){
    len[0] = 0;
    for (int i = 1; i < sz; i++){
        if (f_pos[i] + len[i] - 1 >= s_size){
            len[i] = (s_size - f_pos[i]);
        }
    }
}
void print_suffix_tree(int from){
    cout << "Edge entering in " << from << " has size " << len[from];
    cout << " and starts in " << f_pos[from] << endl;
    cout << "Node " << from << " goes to: ";
    for (auto u : to[from]){
        cout << u.se << " with " << (char)u.fi << " ";
    }
    cout << endl;
    for (auto u : to[from]){
        print_suffix_tree(u.se);
    }
}
void build(string &s){
    for (int i = 0; i < sz; i++){
        to[i].clear();
    }
    sz = 1;
    node = pos = n = 0;
    len[0] = inf;
    for(int i = 0; i < sz(s); i++)
        add_letter(s[i]);
    correct(sz(s));
}
void cutGeneralized(vi &finishPoints){
    for (int i = 0; i < sz; i++){
        int init = f_pos[i];
        int end = f_pos[i] + len[i] - 1;
        int idx = lower_bound(all(finishPoints), init) - finishPoints.begin();
        if ((idx != sz(finishPoints)) && (finishPoints[idx] <= end)){//Must be cut
            len[i] = (finishPoints[idx] - f_pos[i] + 1);
            to[i].clear();
        }
    }
}
void build_generalized(vector<string> &ss){
    for (int i = 0; i < sz; i++){
        to[i].clear();
    }
    sz = 1;
    node = pos = n = 0;
    len[0] = inf;
    int sep = 256;
    vi finishPoints;
    int next = 0;
    for (int i = 0; i < sz(ss); i++){
        for (int j = 0; j < sz(ss[i]); j++){
            add_letter(ss[i][j]);
        }
        next += sz(ss[i]);
        finishPoints.pb(next);
        add_letter(sep++);
        next++;
    }
    correct(next);
    cutGeneralized(finishPoints);
}
\end{minted}